\documentclass{BeamerTemplate}
\title{Chapitre 4 - Équations d'estimation généralisées (GEE)}
\subtitle{STT-4300}
\date{\today}

\IfFileExists{../STT-4300-references.bib}{%
	\addbibresource{../STT-4300-references.bib}%
}{%
	\addbibresource{STT-4300/STT-4300-references.bib}%
}

\begin{document}
\begin{refsection}

\begin{frame}[t]
\titlepage
\end{frame}

\section{Équations d'estimation généralisées (GEE)}
\label{GEE.section}

\begin{frame}{Données longitudinales}
Dans les {\color{blue}études longitudinales} (\textit{longitudinal studies}) (par opposition
aux études transversales), les individus sont mesurés de façon répétée au fil du
temps. Les données longitudinales sont aussi appelées {\color{blue}données de panel} (\textit{panel data}) (surtout en
économie).

\bigskip
L'indépendance est supposée entre les individus (comme dans les MLG), mais pas entre
les mesures d'un même individu.

\bigskip
Les études longitudinales permettent de distinguer les changements au fil du temps
au sein des individus.

\bigskip
Parfois, des {\color{blue}grappes} (\textit{clusters}) sont définies par d'autres groupes qu'un
individu (pratiques médicales, villes, etc.).

\bigskip
\textbf{Notation} : voir page~\pageref{Ximatrix}.
\end{frame}

\begin{frame}{Approches possibles pour les mesures répétées}
Modéliser $Y_{it}$ au moyen de $x_{it}$ :
\begin{itemize}
\item {\color{blue}Modèles marginaux} (\textit{marginal models}) (comme dans les études transversales) :\\
non seulement $g(E(Y_{it})) = x_{it}^T \beta$ ($t=1,\ldots,n_i$), mais aussi $Var(Y_i) = V_i(\phi,\alpha)$.
\item {\color{blue}Modèle à effets aléatoires (ou mixtes)} (\textit{random/mixed effects model}) : la corrélation entre
  les réponses répétées est
  impliquée par la variation des coefficients de régression entre les individus :\\
$g(E(Y_{it}|\gamma_i)) = x_{it}^T \beta + z_{it}^T \gamma_i$, et $\gamma_i$ sont iid
d'une distribution avec $E(\gamma_i)=0$ et $Var(\gamma_i)= \sigma^2_{\gamma}I$.

Une vaste littérature existe pour ces modèles dans le cas gaussien avec lien identité,
beaucoup moins pour d'autres distributions de la famille exponentielle (modèles GLMM).
\end{itemize}

\end{frame}

\subsection{Spécification du modèle}
\begin{frame}[allowframebreaks]{Modèles marginaux}
L'inférence tirée du modèle marginal est de type {\color{blue}moyenne populationnelle} (\textit{population average}) par opposition à l'inférence {\color{blue}spécifique au sujet} (\textit{subject specific}) que l'on peut tirer des modèles à effets aléatoires.

\bigskip
Ingrédients du modèle :
\begin{itemize}
\item L'espérance marginale de la réponse $E(Y_{it}) = \mu_{it}$
  dépend d'un ensemble de variables explicatives $x_{it}$ via $g(\mu_{it})
  = x_{it}^T \beta$, où $g$ est la fonction de lien (mêmes choix ou choix similaires à ceux des MLG).
\item La variance marginale dépend de la moyenne marginale $Var(Y_{it})
  = \phi v(\mu_{it}) $.
\item La corrélation entre $Y_{it}$ et $Y_{it^\prime}$ est une fonction des
  moyennes marginales et possiblement de paramètres supplémentaires $\alpha$.
\end{itemize}

\bigskip
$\beta$ est le paramètre d'intérêt, $\phi$ et $\alpha$ sont des paramètres de nuisance.

\bigskip
Les paramètres $\beta$ s'interprètent comme pour les MLG.

\bigskip
Une distribution marginale pour $Y_{it}$ est postulée (p.~ex.\ Bernoulli, Poisson). Notons,
cependant, que cela ne définit pas une distribution conjointe multivariée pour $Y_i$
($\leadsto$ pas de vraisemblance).

\bigskip\bigskip\bigskip\bigskip\bigskip\bigskip\bigskip
Les paramètres de régression $\beta$ sont estimés par l'approche des
{\color{blue}équations d'estimation généralisées} (\textit{generalized estimating equations}, GEE).
Par analogie avec les équations de score d'un MLG (voir
p.~\pageref{EE.GLM}), on résout (pour $\beta$) :
$$\sum_{i=1}^n D_i^T V_i(\alpha)^{-1} (Y_i - \mu_i) =0,$$
avec $D_i = d \mu_i/ d \beta$, $V_i(\alpha) = \phi  A_i^{1/2} R(\alpha)
A_i^{1/2}$, où $A_i = \mbox{diag}(v(\mu_{it}))$ et
$R(\alpha)$ est une matrice de {\color{blue}corrélation « de travail »} (\textit{working correlation matrix}) (par
opposition à la vraie matrice de corrélation $Corr(Y_i)$).

\bigskip
Les GEE peuvent s'écrire sous la forme d'un algorithme itératif des moindres carrés repondérés.
\end{frame}

\begin{frame}{Choix pour la matrice de corrélation de travail $R(\alpha)$}
\begin{description}
\item[{\color{blue}indépendance}] $R(\alpha)=I$, et on a en fait un
  modèle MLG.
\item[{\color{blue}fixe}] $R(\alpha)$ a une forme prédéfinie connue.
\item[{\color{blue}échangeable} (\textit{exchangeable})] toutes les corrélations $R(\alpha)_{t,t^\prime}$
sont égales à $\alpha$.
\item[{\color{blue}autorégressive (AR)} (\textit{autoregressive})] la corrélation décroît avec
  la différence temporelle, p.~ex.\ $R(\alpha)_{t,t^\prime} = \alpha^{|t-t^\prime|}$.
\item[{\color{blue}$m$-dépendance} (\textit{$m$-dependence})] Les observations sont corrélées jusqu'à
  une distance temporelle $m$.
\item[{\color{blue}non structurée} (\textit{unstructured})] $R(\alpha)$ est
  complètement libre (outre une diagonale de 1 et la contrainte de symétrie).
\end{description}
\end{frame}

\begin{frame}{Estimateurs de $\phi$ et $\alpha$}

Une procédure qui itère entre un
scoring de Fisher modifié pour $\beta$ et une estimation (par les moments) de
$\alpha$ et $\phi$ est utilisée.

\bigskip
Si on définit $r_{it} = (y_{it} -
\mu_{it})/\sqrt{v(\mu_{it})}$, on a
$Var(r_{it}) = \phi$. Un estimateur de $\phi$ peut être
obtenu par ( $N=\sum_{i=1}^n n_i$)
$$\hat{\phi} = \sum_{i=1}^n \sum_{t=1}^{n_i}
\frac{\hat{r}_{it}^2}{N-(p+1)}.$$

\smallskip
L'estimateur spécifique de $R(\alpha)$ dépend du choix de la
structure de corrélation. L'approche générale consiste à estimer $\alpha$ par
une fonction simple de couples de résidus $\hat{r}_{iu}, \hat{r}_{iv}$.

\smallskip
Certaines structures de corrélation exigent des grappes de même taille, c'est-à-dire
$n_i = n_{same}$ pour tout $i$.
\end{frame}

\begin{frame}[allowframebreaks]{Exemples}



\begin{itemize}
\item Si $R(\alpha)_{t,t^\prime}=\alpha$ (échangeable) pour tout $t \neq
  t^\prime$, alors étant donné $\hat{\phi}$, on a
$$\hat{\alpha} = \frac{1}{\hat{\phi}} \sum_{i=1}^n \sum_{t>t^\prime} \hat{r}_{it}
\hat{r}_{it^\prime} /(K-(p+1)),$$
où $K = \sum_{i=1}^n 1/2 \ n_i(n_i-1)$.
\item Si $(R_{\alpha,i})_{tt^\prime} = \alpha^{|t-t^\prime|}$
  (autorégressive), une option est d'utiliser
$$\hat{\alpha}_{t,t^\prime} = \sum_{i=1}^n \frac{\sum_{t=1}^{n_i-(t-t^\prime)}\hat{r}_{it}
\hat{r}_{it^\prime}}{n_i}.$$
\item $\alpha = (\alpha_1, \ldots, \alpha_{n_{same}-1})$, où
  $\alpha_t = R(\alpha)_{t,t+1}$. Alors
$$\hat{\alpha}_t = \frac{1}{\hat{\phi}} \sum_{i=1}^n \hat{r}_{it} \hat{r}_{i(t+1)} /
(n-(p+1)).$$


Si $R(\alpha)$ est tridiagonale avec $R(\alpha)_{t,t+1}=\alpha_t$
(modèle à 1-dépendance), alors si on pose $\alpha_t=\alpha$, on peut
l'estimer par
$$\hat{\alpha} =  \sum_{t=1}^{n_{same}-1} \hat{\alpha}_t /(n_{same}-1).$$

\bigskip
Extension à la $m$-dépendance possible.
\item Si $R(\alpha)$ est totalement non spécifiée, on utilise
$$\hat{R} = \frac{1}{\hat{\phi}n} \sum_{i=1}^n \hat{A}_i^{-1/2} (Y_i - \hat{\mu}_i)
(Y_i - \hat{\mu}_i)^T \hat{A}_i^{-1/2}.$$
\end{itemize}
\end{frame}

\begin{frame}{Remarques}
\begin{itemize}
\item Pour les structures de corrélation d'indépendance, échangeable et de $m$-dépendance,
  $\phi$ n'a pas besoin d'être calculé pour résoudre
  les équations d'estimation (il s'annule). Il est nécessaire pour la structure
  de corrélation AR.
\item Le choix de corrélation échangeable permet des nombres différents
  d'observations et de temps d'observation.
\item La corrélation AR peut accommoder un nombre et un espacement
  arbitraires d'observations.
\end{itemize}

\bigskip

Extension des GEE pour inclure des effets aléatoires :
{\color{red} $\bigstar$ Modèles mixtes pour données non normales, voir
\citeN{Zege:Lian:Albe:1988}.}
\end{frame}


\begin{frame}{Distribution asymptotique de l'estimateur GEE}
Si un estimateur $\sqrt{n}$-convergent est utilisé pour estimer $\alpha$ et
$\phi$, on
peut montrer que $\sqrt{n} (\hat{\beta} - \beta)$ est asymptotiquement
normal avec variance
$$ \Omega = \lim_{n \rightarrow \infty} n M^{-1} Q M^{-1},$$
où
$$M = \sum_{i=1}^n D_i^T V_i^{-1} D_i,$$
et
$$Q = \sum_{i=1}^n D_i^T V_i^{-1} Var(Y_i) V_i^{-1} D_i.$$
\end{frame}

\begin{frame}{Remarques}
\begin{itemize}
\item La variance asymptotique de $\hat{\beta}$ ne dépend pas du
  choix des estimateurs de $\alpha$ et $\phi$ parmi ceux
  qui sont $\sqrt{n}$-convergents.
\item La convergence de $\hat{\beta}$ (et $\hat{\Omega}$, voir page~\ref{Omega.page}) ne dépend
que de la spécification correcte
de la moyenne $\mu_i$ et non de la spécification correcte
de la structure de corrélation.
\item Choisir $R$ proche de la vraie matrice de corrélation augmente
  l'efficacité.
\item L'inférence et les diagnostics sont limités.
\end{itemize}
\end{frame}

\begin{frame}{Diagnostics}

Comme pour les MLG, on définit les résidus de Pearson (\textit{Pearson residuals})
$$\hat{r}_{it} = \frac{y_{it} - \hat{\mu}_{it}}{\sqrt{v(\hat{\mu}_{it})}}.$$

\bigskip
Ils peuvent être représentés graphiquement pour identifier les valeurs aberrantes et autres violations des hypothèses.

\end{frame}

\begin{frame}{Inférence}
On utilise la théorie asymptotique pour dériver des intervalles de confiance
et des tests z (approximatifs).

\label{Omega.page}
\bigskip
Un estimateur de la variance asymptotique $\Omega$ est nécessaire. Celui-ci est
$\hat{\Omega} = \hat{M}^{-1} \hat{Q} \hat{M}^{-1},$
où
$$\hat{M} = \sum_{i=1}^n \hat{D}_i^T \hat{V}_i^{-1} \hat{D}_i,$$
et
$$\hat{Q} = \sum_{i=1}^n \hat{D}_i^T \hat{V}_i^{-1} (Y_i - \hat{\mu_i}) (Y_i -
\hat{\mu_i})^T \hat{V}_i^{-1} \hat{D}_i.$$

\smallskip
$\hat{D}_i$ et $\hat{V}_i$ sont obtenus en substituant $\hat{\beta}$,
$\hat{\phi}$ et $\hat{\alpha}$.
\end{frame}

\begin{frame}{}
L'absence d'une fonction de vraisemblance pour ces modèles rend l'inférence
limitée.

\bigskip
\bigskip
{\color{red} $\bigstar$ Un critère de type $C_p$ pour la sélection de variables
  dans les modèles longitudinaux marginaux, voir \citeN{Cant:Flem:Ronc:2005}.}

\bigskip
{\color{red} $\bigstar$ Un critère de type Akaike pour les GEE, voir
  \citeN{Pan:2001}.}
\end{frame}

\subsection{Jeu de données CARDIA}
\begin{frame}[fragile]{Analyse du jeu de données CARDIA (voir p.~\pageref{CARDIApage})}
{\tiny
\begin{verbatim}
Call: gee(formula = smoke ~ age + factor(birth) + factor(education) + 
    factor(racesex), id = id, data = CARDIA.sub, family = binomial, 
    corstr = "exchangeable")

Coefficients:
                      Estimate Naive S.E.    Naive z Robust S.E.   Robust z
(Intercept)        -1.33707777 0.54231807  -2.465486  0.54510289  -2.452891
age                 0.06480087 0.02646754   2.448315  0.02669744   2.427231
factor(birth)2     -0.19856466 0.14606724  -1.359406  0.14563751  -1.363417
factor(birth)3     -0.40380459 0.23689643  -1.704562  0.23858001  -1.692533
factor(education)2 -0.68712005 0.08313860  -8.264754  0.08490269  -8.093030
factor(education)3 -1.96140806 0.10016299 -19.582164  0.10054908 -19.506972
factor(racesex)2   -0.17668132 0.09763816  -1.809552  0.09900202  -1.784623
factor(racesex)3   -0.13445604 0.10634023  -1.264395  0.10833799  -1.241079
factor(racesex)4   -0.12730906 0.10530815  -1.208919  0.10599195  -1.201120

Estimated Scale Parameter:  1.000362
Number of Iterations:  1

Working Correlation
         [,1]     [,2]     [,3]     [,4]
[1,] 1.000000 0.727139 0.727139 0.727139
[2,] 0.727139 1.000000 0.727139 0.727139
[3,] 0.727139 0.727139 1.000000 0.727139
[4,] 0.727139 0.727139 0.727139 1.000000
\end{verbatim}
}
\end{frame}

\begin{frame}[fragile]{}
{\tiny
\begin{verbatim}
Call: gee(formula = smoke ~ age + factor(birth) + factor(education) + 
    factor(racesex), id = id, data = CARDIA.sub, family = binomial, 
    corstr = "AR-M", Mv = 1)

Coefficients:
                      Estimate Naive S.E.    Naive z Robust S.E.   Robust z
(Intercept)        -1.46502183 0.51848241  -2.825596  0.54981825  -2.664557
age                 0.07058083 0.02530312   2.789412  0.02691652   2.622212
factor(birth)2     -0.21226645 0.13970893  -1.519348  0.14716072  -1.442412
factor(birth)3     -0.41934193 0.22641663  -1.852081  0.24111643  -1.739168
factor(education)2 -0.70331577 0.07936239  -8.862080  0.08536346  -8.239072
factor(education)3 -1.98023517 0.09576186 -20.678745  0.10143031 -19.523112
factor(racesex)2   -0.19210488 0.09334897  -2.057922  0.09978203  -1.925245
factor(racesex)3   -0.14273599 0.10162928  -1.404477  0.10912066  -1.308056
factor(racesex)4   -0.11428211 0.10047526  -1.137415  0.10674042  -1.070654

Estimated Scale Parameter:  1.006728
Number of Iterations:  2

Working Correlation
          [,1]      [,2]      [,3]      [,4]
[1,] 1.0000000 0.7684591 0.5905294 0.4537977
[2,] 0.7684591 1.0000000 0.7684591 0.5905294
[3,] 0.5905294 0.7684591 1.0000000 0.7684591
[4,] 0.4537977 0.5905294 0.7684591 1.0000000
\end{verbatim}
}
\end{frame}

\begin{frame}[fragile]{}
{\tiny
\begin{verbatim}
Call: gee(formula = smoke ~ age + factor(education) + factor(racesex), 
    id = id, data = CARDIA.sub, family = binomial, corstr = "exchangeable")

Coefficients:
                      Estimate  Naive S.E.    Naive z Robust S.E.   Robust z
(Intercept)        -0.51416495 0.251656751  -2.043120 0.253634983  -2.027185
age                 0.02272524 0.009743051   2.332456 0.009868721   2.302754
factor(education)2 -0.68987727 0.082992723  -8.312503 0.084840099  -8.131500
factor(education)3 -1.96816994 0.100008335 -19.680059 0.100455198 -19.592515
factor(racesex)2   -0.17769210 0.097552550  -1.821501 0.098969587  -1.795421
factor(racesex)3   -0.13115039 0.106197596  -1.234966 0.108213885  -1.211955
factor(racesex)4   -0.12631590 0.105246698  -1.200189 0.106046722  -1.191134

Estimated Scale Parameter:  0.9997882
Number of Iterations:  1

Working Correlation
          [,1]      [,2]      [,3]      [,4]
[1,] 1.0000000 0.7272402 0.7272402 0.7272402
[2,] 0.7272402 1.0000000 0.7272402 0.7272402
[3,] 0.7272402 0.7272402 1.0000000 0.7272402
[4,] 0.7272402 0.7272402 0.7272402 1.0000000
\end{verbatim}
}
\end{frame}

\begin{frame}[fragile]{}
{\tiny
\begin{verbatim}
Call: gee(formula = smoke ~ age + factor(education), id = id, data = CARDIA.sub, 
    family = binomial, corstr = "exchangeable")

Coefficients:
                      Estimate  Naive S.E.    Naive z Robust S.E.   Robust z
(Intercept)        -0.59982654 0.247153965  -2.426935 0.248914981  -2.409765
age                 0.02204566 0.009672377   2.279239 0.009739698   2.263485
factor(education)2 -0.70520857 0.082571026  -8.540630 0.084484169  -8.347227
factor(education)3 -1.98709483 0.096223429 -20.650842 0.096193423 -20.657284

Estimated Scale Parameter:  0.9995199
Number of Iterations:  1

Working Correlation
          [,1]      [,2]      [,3]      [,4]
[1,] 1.0000000 0.7275345 0.7275345 0.7275345
[2,] 0.7275345 1.0000000 0.7275345 0.7275345
[3,] 0.7275345 0.7275345 1.0000000 0.7275345
[4,] 0.7275345 0.7275345 0.7275345 1.0000000
\end{verbatim}
}
\end{frame}

\begin{frame}{}
\begin{center}
\includegraphics[width = 7cm]{figures/CARDIA3resid.pdf}  
\end{center}
\end{frame}

\subsection{Jeu de données Epileptic}
\begin{frame}{Analyse du jeu de données Epileptic}
Essai clinique avec 59 patients épileptiques, randomisés soit au médicament
antiépileptique progabide, soit au placebo (tiré de \citeNP{Thal:Vail:1990}).

\bigskip
Avant le traitement, le nombre de crises dans une période de huit semaines
a été compté pour chaque patient comme mesure de référence
(\textit{baseline}). Le nouveau traitement est-il efficace pour réduire le taux de crises chez les patients épileptiques?

\bigskip
On commence avec un modèle de Poisson, mais... surdispersion possible...

\bigskip
  \begin{tabular}{lcccc}
\hline
\multicolumn{1}{l}{} & \multicolumn{3}{c}{Visite}\\
\hline
Traitement & 1 & 2 & 3 & 4\\
\hline
Traité & 38.7 & 16.8 & 23.8 & 18.8 \\
Placebo & 10.8 & 7.5 & 24.5 & 7.3\\
\hline
  \end{tabular}
\end{frame}


\begin{frame}{}
De plus, les temps d'observation sont : après 8 (référence), 10, 12, 14 et 16
semaines. Il faut tenir compte de la différence de durée de la
période d'observation. On utilise un décalage (\textit{offset}) de la forme $\log(t_{it})$,
où $t_{it}=8$ pour la première visite et $t_{it}=2$ pour les
suivantes.

\bigskip
Une variable \texttt{post} pour tenir compte du début de l'étude,
qui prend la valeur 1 pour les visites 1, 2, 3 et 4 et la valeur 0 pour
la visite de référence.

\bigskip
Permettre une interaction entre \texttt{post} et le groupe.

\end{frame}

\begin{frame}[fragile]{}
{\tiny
\begin{verbatim}
Call: gee(formula = y ~ post * factor(group) + offset(log(rep(c(8, 
    2, 2, 2, 2), 59))), id = Subject, data = seizure, family = poisson, 
    corstr = "exchangeable")

Coefficients:
                               Estimate Naive S.E.    Naive z Robust S.E.   Robust z
(Intercept)                  1.34760922  0.1510969  8.9188397   0.1573571  8.5640166
post                         0.11183602  0.1545145  0.7237900   0.1159304  0.9646821
factor(group)progabide       0.02753449  0.2071018  0.1329515   0.2217878  0.1241479
post:factor(group)progabide -0.10472579  0.2197052 -0.4766650   0.2134448 -0.4906459

Estimated Scale Parameter:  19.6797
Number of Iterations:  1

Working Correlation
          [,1]      [,2]      [,3]      [,4]      [,5]
[1,] 1.0000000 0.7713861 0.7713861 0.7713861 0.7713861
[2,] 0.7713861 1.0000000 0.7713861 0.7713861 0.7713861
[3,] 0.7713861 0.7713861 1.0000000 0.7713861 0.7713861
[4,] 0.7713861 0.7713861 0.7713861 1.0000000 0.7713861
[5,] 0.7713861 0.7713861 0.7713861 0.7713861 1.0000000
\end{verbatim}
}
\end{frame}

\begin{frame}{}
\begin{center}
\includegraphics[width = 6.5cm]{figures/seizureresid.pdf}  
\end{center}
Une très grande valeur aberrante!
\end{frame}

\begin{frame}[fragile]{Sans le patient 49}
{\tiny
\begin{verbatim}
Call: gee(formula = y ~ post * factor(group) + offset(log(rep(c(8, 
    2, 2, 2, 2), 59))), id = Subject, data = seizure, subset = Subject != 
    49, family = poisson, corstr = "exchangeable")

Coefficients:
                              Estimate Naive S.E.    Naive z Robust S.E.   Robust z
(Intercept)                  1.3476092  0.1105249 12.1928162   0.1573571  8.5640166
post                         0.1118360  0.1231346  0.9082421   0.1159304  0.9646821
factor(group)progabide      -0.1068224  0.1578051 -0.6769263   0.1936977 -0.5514904
post:factor(group)progabide -0.3023841  0.1933863 -1.5636272   0.1710601 -1.7677071

Estimated Scale Parameter:  10.52997
Number of Iterations:  1

Working Correlation
         [,1]     [,2]     [,3]     [,4]     [,5]
[1,] 1.000000 0.593689 0.593689 0.593689 0.593689
[2,] 0.593689 1.000000 0.593689 0.593689 0.593689
[3,] 0.593689 0.593689 1.000000 0.593689 0.593689
[4,] 0.593689 0.593689 0.593689 1.000000 0.593689
[5,] 0.593689 0.593689 0.593689 0.593689 1.000000
\end{verbatim}
}

Grandes différences dans (certains) paramètres estimés et erreurs
types!
\end{frame}

\begin{frame}{Remerciements}
Ces diapositives ont été inspirées par le cours GLAM de Eva Cantoni de l'Université de Genève.
\end{frame}

\begin{frame}[allowframebreaks]{Références du chapitre}
\tiny
\printbibliography[heading=none,section=\therefsection]
\end{frame}

\end{refsection}

\end{document}
